Fix a realized graph \(A\) and let \(S_0=\{i:N_i=0\}\).  Consider the single admissible response law
\[
 f_i(a,x)=\varepsilon_iBx,
 \qquad \varepsilon_i\overset{\mathrm{iid}}\sim{\rm Unif}\{-1,1\},
\]
independent of the graph variables.  It obeys the response envelope because \(|f_i(a,x)|\leq B\), \(|\partial_xf_i(a,x)|=B\), and its second and third derivatives vanish; thus it also obeys the third-derivative bound for every \(L\geq0\).  Its use does not alter any graph, rank, support, or randomization property of \({\rm def:model-class}\).

Condition on two response-sign configurations that agree on \(S_0^c\) and are opposite on \(S_0\).  Both configurations have positive conditional probability under this one i.i.d. response law.  For every isolate, anonymous interference gives
\[
 Y_i^{\mathrm{obs}}=f_i(Z_i,0)=0,
\]
and an isolate's response sign enters no other unit's response.  Hence the complete conditional observed-data laws are identical under the two configurations.  Their targets, however, differ by
\[
 \frac{2B}{n}|S_0|=\frac{2B}{n}I_0(A).
\]
If \(m(A)\) denotes the common conditional mean of an arbitrary estimator under these indistinguishable configurations and \(\theta^+,\theta^-\) their targets, the triangle inequality gives
\[
 \max\{|m(A)-\theta^+|,|m(A)-\theta^-|\}
 \geq\frac12|\theta^+-\theta^-|
 =\frac{B}{n}I_0(A).
\]
Therefore any common graph-only conditional-bias radius valid for every response configuration must be at least this large.  The argument is pointwise in \(A\), so averaging cannot turn this conditional radius into the smaller centered fluctuation scale.